\chapter{Perturbations about a sewn solution}\label{ch:12}
Send in two weak pulses. A linear field gives the sum of the two signals,
but an interacting field can respond to their overlap. The first coefficient
sensitive to that overlap is a mixed derivative with respect to the source
strengths. Its calculation makes a good test of sewing: the pulses, the
vertex where they meet, and the detector need not be on the same side. We
will compute the coefficient twice, first from a classical initial-value
problem and then as a tree term of finite-time quantum evolution.

\section{Choosing the background and the free problem}
Let $\bar\phi$ be a solution on the whole slab, or a matched regional pair
proved to assemble to such a solution. Fix the outer data and geometric
sector. Write $\phi=\bar\phi+\eta$ in a field chart. A source expansion
will vary the bulk driving force while holding initial data fixed. Thus the
response $\eta$ has zero incoming Cauchy data, although the background
need not have zero data. With the exterior conditions held fixed, and for
variations supported away from the caps, the first variation of the joined
action vanishes on $\bar\phi$. This includes the cancellation of its two
outward responses. Otherwise the proposed fluctuation theory would start
with a source on the seam even before the pulses are sent.

The quadratic action is one half of the second variation evaluated twice
on $\eta$; its operator is the Jacobi operator at $\bar\phi$. Its domain includes the linearized transmission and
outer conditions. Higher variations give interaction vertices, including
any allowed physical boundary vertices. Linearizing independently around
two convenient regional vacua need not give these operators: their fields
or normal momenta might not be the traces of one solution. For a gauge
theory a Green operator may require a gauge choice; the external
observables must still descend to the physical theory of
Chapter~\ref{ch:11}. We now choose a scalar self-interaction for which the
mixed response can be found without any gauge assumption.

Use a fixed background metric, $P=\Box-m^2-\xi R$, and add to the free
action the terms
\begin{equation}
 -\int_M\left(\frac{\lambda}{3!}\chi\phi^3+j\phi\right)dV .
 \label{eq:12-source-action}
\end{equation}
Here $\chi$ is smooth, $j$ is an applied bulk source, and $\lambda$ counts
cubic vertices. Use compact switching where end effects or spatial support
matter. The action sign fixes the equation
$P\phi-\lambda\chi\phi^2/2=j$. Suppose $\bar\phi$ solves it at $j=j_0$,
with the transparent scalar matching of Chapter~\ref{ch:04}. For
$j=j_0+s f+r h$ the \emph{exact} equation for the fluctuation is
\begin{equation}
 P_{\bar\phi}\eta=s f+r h+\tfrac{\lambda}{2}\chi\eta^2,
 \qquad P_{\bar\phi}=P-\lambda\chi\bar\phi .
 \label{eq:12-jacobi}
\end{equation}
The term linear in $\eta$ from expanding the interaction belongs to the
free operator about this solution. The quadratic coefficient is
$\tfrac12\int\eta P_{\bar\phi}\eta\,dV$ when the integrations by parts
have no uncancelled surface terms; the remaining interaction is cubic in
$\eta$. Since the interaction in \eqref{eq:12-source-action} contains no
derivatives, it does not change the scalar trace or the outward normal
momentum. This is one reason for the choice of model.

\section{The mixed response and its regional construction}
Let $G_R^{\bar\phi}$ be the retarded inverse of $P_{\bar\phi}$ on this
problem. Assume the finite-time initial-boundary problem is well posed for
the smooth sources used here and has the requisite derivatives at zero;
alternatively read the following as a formal recursion in $s,r$ with a
well-defined linear inverse. Write $u_f=\partial_s\eta|_0$,
$u_h=\partial_r\eta|_0$, $w_{fh}=\partial_s\partial_r\eta|_0$.
Differentiating \eqref{eq:12-jacobi} and imposing zero initial data at each
order gives
\begin{align}
 u_f&=G_R^{\bar\phi} f,\qquad u_h=G_R^{\bar\phi}h,\nonumber\\
 P_{\bar\phi}w_{fh}&=\lambda\chi u_f u_h,\qquad
 w_{fh}=\lambda G_R^{\bar\phi}\!\left(\chi u_f u_h\right).
 \label{eq:12-response}
\end{align}
There are two ways to assign the incoming derivatives to the two factors
of $\eta^2$. They cancel the $1/2$ in its coefficient. Thus the coefficient
of $sr$ in $\eta$ is $w_{fh}$, whereas the coefficient of $s^2$ alone is
$\lambda G_R^{\bar\phi}(\chi u_f^2)/2$. The expression retains the
background dependence of the inverse; it is not an assertion that the
background is of order zero in an independent expansion in $\lambda$.

For a numerical sign and normalization check, take a massless flat patch,
$\bar\phi=0$, and $\chi=f=h=1$ in the causal past of the observation
between the initial cap and elapsed time $T$. Cutoffs outside that domain
do not change this retarded calculation. With zero initial perturbation
the solution there has no spatial dependence. Its first derivative is
$u_f(T)=u_h(T)=-T^2/2$, because $P=-\partial_T^2$ on this field. The
second derivative is
\[
 w_{fh}(T)=-\frac{\lambda}{4}\int_0^T(T-t)t^4\,dt
          =-\frac{\lambda T^6}{120} .
\]
In particular the nonlinear signal is not obtained by simply multiplying
the two first-order readings at the detector. Their product acts as a new
source; it must propagate from its point of formation to the detector.

To reproduce \eqref{eq:12-response} from regional data, first compute
$u_f,u_h$ by solving the linearized response problem on the pieces. On a
timelike scalar cut this means varying their common boundary history until
the two outward normal responses cancel. Do the same for the source
$\lambda\chi u_f u_h$ at the next step. Why should the result equal a
whole-slab response? For this scalar, continuity kills a delta derivative
in the wave operator of a piecewise field; equality of the two normal
derivatives in a common signed coordinate kills the delta itself. The
piecewise field then solves the whole inhomogeneous equation weakly. With
the prescribed past and outer data, uniqueness in the solution class gives
the claimed response. This uses the reconstruction and Green-operator
hypotheses of Chapters~\ref{ch:07} and~\ref{ch:08}; it is not a rule for
manipulating Heaviside functions times quantum fields.

There is a useful way to compare the two computations directly without
placing a singular Green kernel on the cut. Smear the output with a smooth
detector profile $k$. By adjointness of retarded and advanced inverses,
\begin{equation}
 \int_M k\,w_{fh}\,dV
 =\lambda\sum_{a=1,2}\int_{M_a}
 (G_A^{\bar\phi}k)\,\chi\,u_f u_h\,dV .
 \label{eq:12-sewn-tree}
\end{equation}
The restriction to each piece is now an integral of smooth fields. The
advanced solution carries the detector profile back to the vertex; the
retarded solutions carry the sources forward. If $f$ and $h$ lie in
different pieces, replacing the transmitted fields by independently
reflecting fields can set the overlap to zero. With the correctly joined
Green operators the two integrals instead add to the whole answer, without
an independent interaction assigned to the artificial wall. Were the
interface equation itself nonlinear, differentiating it would produce an
additional inhomogeneous interface condition for $w_{fh}$; it is absent
here by the particular action and domain.

\section{Finite-time evolution and the same tree}
Choose the same free fluctuation problem, with an admissible state if an
expectation is wanted. Write $\widehat\eta_I$ for the interaction-picture
free field. For this nondifferentiated potential the interaction Hamiltonian
has the opposite sign to the added Lagrangian density. For a foliation with
$dV=N\,dt\,d\Sigma_t$ its insertions are
\[
 V_I(t)=\int_{\Sigma_t} N\,d\Sigma_t\left[
  \frac{\lambda}{3!}\chi\, :\widehat\eta_I^3:
     +(s f+r h)\widehat\eta_I\right].
\]
The colons denote Wick products in a specified subtraction. In the field
theory $V_I(t)$ denotes an operator-valued distribution in time obtained
from spacetime insertions, integrated with smooth switching; the notation
does not assert that a sharp-time Wick cube is an operator on a Hilbert
space. No stationary Hamiltonian or choice of particles at infinity is
needed: the free linear evolution identifies the slice algebras. With
$\hbar=1$, the Dyson series on the finite slab is
\begin{align}
 U(t_f,t_i)&=\mathcal T\exp\!\left(-i\int_{t_i}^{t_f} V_I(t)\,dt\right)
 \nonumber\\
 &=\mathbf 1-i\int_{t_i}^{t_f}V_I(t)\,dt
 -\frac{1}{2!}\int_{t_i}^{t_f}dt\int_{t_i}^{t_f}dt'\,
   \mathcal T\{V_I(t)V_I(t')\}+\cdots .
 \label{eq:12-dyson}
\end{align}
Later times stand on the left. In a field theory this is an order-by-order
formula for smeared, renormalized insertions, not an operator-norm
convergence claim. One may use switching supported in the interior of the
slab. If an actual time-dependent Hamiltonian evolution is defined,
$U(t_f,t_i)=U(t_f,t_*)U(t_*,t_i)$ follows by splitting the ordered
integration domains. For local perturbations the corresponding statement
is causal factorization for suitably ordered supports. Neither operation
makes two spatial regions independent, and neither takes an infinite-time
scattering limit \cite{FR}.

The quantity appropriate to our detector is a response in the specified
initial state, namely the expectation of
$U(t_x,t_i)^*\widehat\eta_I(x)U(t_x,t_i)$, understood after smearing.
At first order in $V_I$, its change is
$i\int_{t_i}^{t_x}[V_I(t),\widehat\eta_I(x)]\,dt$. For the source term,
$i[\widehat\eta_I(y),\widehat\eta_I(x)]=E_{\bar\phi}(x,y)\mathbf1$.
With retarded ordering the result is $u=G_R^{\bar\phi}(s f+r h)$.
All higher source commutators of this c-number vanish: in the free theory
a source simply displaces the field by $u$. To keep that displacement
exact while counting cubic vertices, factor out the evolution $U_j$
generated by the source alone, writing $U=U_j U_{\rm rel}$. The cubic
insertion in $U_{\rm rel}$ is then conjugated by $U_j$: its field at $y$
is $\widehat\eta_I(y)+u(y)$. Use one Wick prescription for the powers,
so that the ordinary binomial rule holds for a shift by a smooth function.
The commutator needed at first cubic order is
\[
 i[:\!(\widehat\eta_I(y)+u(y))^3\!:,\widehat\eta_I(x)]
 =3 E_{\bar\phi}(x,y):\!(\widehat\eta_I(y)+u(y))^2\!: .
\]
The three identical terms cancel a factor of three in $3!$. With retarded
ordering the one-vertex correction to the mean field is therefore
\begin{equation}
 \frac{\lambda}{2}\int_M G_R^{\bar\phi}(x,y)\chi(y)
  \left\langle :\!\big(\widehat\eta_I(y)+u(y)\big)^2\!:\right\rangle dV_y .
 \label{eq:12-displaced-tree}
\end{equation}
Equivalently, insert the source-only solution into the quadratic source on
the right side of the renormalized field equation at first vertex order.
For a centered free state, the term quadratic in source amplitudes is just
$u(y)^2$. Its mixed derivative is $2u_fu_h$, so
\eqref{eq:12-displaced-tree} reproduces \eqref{eq:12-response} with both
its sign and its factor of $\lambda$. A local expectation of the Wick
square, if nonzero in the chosen subtraction, is source independent at
this order. It is not part of the two-pulse response.

\section{What survives the comparison at the next order?}
A classical response graph is drawn with its output at the root and a
retarded line for each inversion of $P_{\bar\phi}$. Equation
\eqref{eq:12-sewn-tree} has one three-legged vertex and no closed chain.
Iterating the classical equation gives more vertices but still trees. In a
quantum calculation two fluctuations at vertices can instead be contracted
with each other. Two lines joining the same pair of vertices already make
one loop. The distinction is not the power of $\lambda$: two cubic
vertices can occur both in a tree and in a loop. For example, with $n$
vertices and $I$ internal lines a connected graph has $L=I-n+1$ independent
loops. One way to keep the counting straight is to expand a connected
functional with external classical fields held fixed: internal contractions
carry $\hbar$, action insertions carry $1/\hbar$, and the overall
$\hbar$ multiplying the connected logarithm gives $\hbar^L$. The tree
response is of order $\hbar^0$ in this counting.

In a centered quasifree reference state Wick's rule reduces free-field
expectations in the Dyson expansion to pair contractions
$G_F(x,y)=\langle\mathcal T\widehat\eta_I(x)\widehat\eta_I(y)\rangle$,
without a further factor of $i$, once the coincident products are defined.
For a general initial state higher connected correlations are additional
data. The time-ordered contraction is not the retarded Green operator used
in a classical response. Even away from coincident vertices, an ordinary
quantum correlator computed in independently chosen regional states need
not have the correct cross contraction. At coincidence the product of two
lines in a loop can also require a local extension. Agreement of the
transmitted Hessian and all the classical vertices gives the tree response
derived here; agreement of a quantum coefficient further requires the
chosen state's contractions, the same operator and subtraction
normalizations, and compatible time-ordered products including contact
terms. Chapter~\ref{ch:13} works out the first example where a genuine
extension choice enters. The one-vertex calculation supplies a comparison
at that order, not an all-orders commutation of quantization with sewing.

\begin{exercise}
Couple a third source $b\,l$ to the same cubic problem. Derive the fully
mixed third derivative of the retarded classical field at zero source in
terms of the first responses and $w_{fh}$ of \eqref{eq:12-response}. How
many distinct terms remain when the three source labels are different?
\end{exercise}
\noindent\emph{Hint.} Differentiate $\eta^2$ three times at $\eta=0$.
One factor must receive one derivative, the other two. Pair the two possible
orders of these factors before applying $G_R^{\bar\phi}$.
